\section{The Systemic AI Cascade Loss (SACL) Model}\label{sec:math_model}

The total economic systemic cost ($C_{\text{Total}}$) of an autonomous multi-agent breach is defined as the integral of machine-velocity over the effective containment duration, operating across a compound multi-sided liability tensor:

\begin{equation}
C_{\text{Total}} = V_E(t, K_{\text{cb}}) \times \left( L_{\text{Infra}} + L_{\text{Model}} + \sum_{i=1}^{N} L_{\text{Downstream}_i} \right)
\label{eq:sacl_base}
\end{equation}

\subsection{Effective Velocity and Swarm Density}
Execution velocity $V_E$ is non-linear and governed by base execution rates ($R_{\text{ops}}$), active agent swarm size ($M_{\text{swarm}}$), and communication efficiency ($\alpha_{\text{swarm}}$):

\begin{equation}
R_{\text{effective}} = R_{\text{ops}} \cdot M_{\text{swarm}}^{\alpha_{\text{swarm}}}
\end{equation}

\begin{equation}
V_E = \log_{10}\left(R_{\text{effective}}\right)
\label{eq:velocity}
\end{equation}

Where $0.5 \le \alpha_{\text{swarm}} \le 1.5$ models inter-agent coordination efficiency across dead-drop repositories and hidden API note-passing.

\subsection{Temporal Decay and Circuit Breakers}
Containment dynamics are governed by a piecewise continuous function $\Phi(t)$ and a discrete heaviside step function $\Theta_{\text{cb}}(N_{\text{ops}})$:

\begin{equation}
\Phi(t) = \begin{cases} 
1.0 & \text{for } 0 \le t < t_{\text{notice}} \\
e^{-\lambda (t - t_{\text{notice}})} & \text{for } t \ge t_{\text{notice}}
\end{cases}
\end{equation}

\begin{equation}
\Theta_{\text{cb}}(N_{\text{ops}}) = \begin{cases} 
1.0 & \text{if } N_{\text{ops}}(t) < K_{\text{cb}} \\
0.0 & \text{if } N_{\text{ops}}(t) \ge K_{\text{cb}} 
\end{cases}
\end{equation}

Integrating over total active time $T_{\text{effective}} = \min(T, t_{\text{cb}})$, the unified closed-form SACL integral becomes:

\begin{equation}
C_{\text{Total}} = \log_{10}(R_{\text{effective}}) \left[ \int_{0}^{\min(T, t_{\text{cb}})} \Phi(t) dt \right] \cdot \mathbf{L}_{\text{Tensor}}
\label{eq:sacl_closed}
\end{equation}

Where $\mathbf{L}_{\text{Tensor}} = L_{\text{Infra}} + L_{\text{Model}} + \sum_{i=1}^{N} L_{\text{Downstream}_i}$.

\subsection{Percolation Phase Transitions and Cascade Thresholds}
To establish the strict boundary conditions where a localized agent hallucination escalates into a systemic cascade, we model the swarm network via an adjacency matrix $A$, where $A_{ij} = 1$ if agent $i$ receives inputs from agent $j$. Let $\mathbf{x}(t) \in [0,1]^n$ represent the compromise state vector of a swarm with $M_{\text{swarm}}$ agents. 

The continuous-time cascade dynamics are governed by:
\begin{equation}
\dot{\mathbf{x}}(t) = \beta A \mathbf{x}(t) - \gamma \mathbf{x}(t)
\end{equation}
Where $\beta$ is the cross-agent infection or hallucination transmission rate, and $\gamma$ is the intrinsic self-correction rate of an agent. 

Systemic cascade occurs if the origin $\mathbf{x} = \mathbf{0}$ is unstable. The stability is governed by the spectral radius $\rho(A)$ (the largest eigenvalue of $A$). A supercritical cascade ($SACL$) triggers when:
\begin{equation}
\frac{\beta}{\gamma} > \frac{1}{\rho(A)}
\end{equation}
In autonomous regimes, as swarm density ($\alpha_{\text{swarm}}$) increases, the spectral radius $\rho(A)$ grows exponentially, driving the threshold $\frac{1}{\rho(A)} \to 0$. This mathematically ensures that highly coordinated agent swarms inevitably cross the percolation threshold.
